\section{Lessons Learned}

\paragraph{A strong project structure and scope is the key.} The complexity of the project grew much more than expected. During a control meeting we realized we were facing real limitations, especially with the AnalystAgent: a single script was not enough to handle the different prompt versions, the section-splitting logic, and the various output formats we wanted to test. We decided to split the AnalystAgent into three separate agents, each writing to its own sentiment table. This decision was crucial, but this solution and the ones implemented to face other challenges would not have been possibile without the clear project structure and scope defined in the planning phase.

\paragraph{Schema first, code second.} The biggest single source of bugs was schema drift between developer machines and the VPS. The team eventually adopted a hard rule, codified in \texttt{AGENTS.md}: \texttt{db/schema.sql} is authoritative; any table or column not declared there must not be read or written; if a feature needs a new column, the schema is updated in the same commit. Every subsequent integration error became easy to diagnose because the failing query named the missing column outright.

\paragraph{Multiple analysts beat one ``perfect'' analyst.} The decision to ship \emph{three} sentiment tables (\texttt{sentiment}, \texttt{sentiment2}, \texttt{sentiment3}) plus a weighted variant (\texttt{sentiment\_w}) felt wasteful at first but turned out to be the cheapest possible ablation framework: each analyst writes to its own table with the same \texttt{document\_id} foreign key, so cross-model comparisons are one SQL join away. We recommend this pattern for any LLM-in-the-loop pipeline where the prompt and segmentation are still being tuned.

\paragraph{A static dashboard scales further than you think.} The dashboard at \url{https://fedwatcher.ellep.it} is hand-written HTML, CSS and one vanilla JavaScript file (\texttt{explorer.js}). It calls one or two FastAPI endpoints and renders everything client-side with Chart.js. No build step, no framework, no node\_modules. Deploying a change is \texttt{scp} of a single file. For a project where the team rotates between four developers and one VPS, this was the right call.

\paragraph{Use AI agents for what they are good at.} Generative AI was used most effectively for (a) writing the AnalystAgent prompt, (b) explaining unfamiliar pieces of the codebase to whichever team member touched them next, and (c) ops scripting (the May~27 database merge was driven by Claude Code through an SSH session). It was used least effectively for designing the schema -- early LLM-suggested tables had inconsistent naming conventions and missing foreign keys, and the schema only stabilized after a human review pass.

\paragraph{Cron + SQLite + systemd is enough.} Several alternatives were considered (Celery + Redis, Airflow, RDS) and all were rejected as over-engineering. Five cron lines and one systemd unit run the entire production pipeline. The full ops surface fits in less than 30 lines of configuration, which made the May~27 server consolidation a one-evening task rather than a sprint.

\paragraph{What we would do differently.} (i) Write the schema and the AGENTS contract on day one, not day forty; (ii) ship the FastAPI service before the static-JSON snapshot, so the dashboard never has to migrate; (iii) integrate the divergence sign convention into a unit test, not a code review; (iv) commit more frequently with smaller diffs -- some of our commits ended up bundling unrelated changes and are harder to attribute in this diary than they should be.
